\documentclass[a4paper,12pt]{article}
\usepackage{color}
\usepackage{graphicx, gensymb}
\usepackage{amsfonts, cancel, amssymb}
\usepackage{amsmath, amsthm, array, esvect, esint}
\usepackage{typearea}
\usepackage[a4paper,left=2cm,right=2cm,top=2cm,bottom=3cm,bindingoffset=5mm]{geometry}
\newcommand\numberthis{\addtocounter{equation}{1}\tag{\theequation}}
\newcommand{\bra}{}
\newcommand{\kom}{}
\newcommand{\brak}{}
\newcommand{\braket}{}
\newcommand{\intx}{}
\newcommand{\intr}{}
\newcommand{\kla}{}
\newcommand{\klam}{}
\newcommand{\klammer}{}
\newcommand{\oper}{}

\begin{document}

\title{05 Beschränkte Operatoren}
\author{Joachim Schwardt}
\date{\today}
\maketitle

\section{Beschränkte lineare Operatoren}
\begin{proof}
Da $A$ beschränkt ist $\exists M\ge 0\forall\phi\in D(A):\|A\phi\|\le M\|\phi\|$. Sei $\phi_n\rightarrow\phi\in\mathcal{H}$ eine Folge in $D(A)$. Dann gilt:
$$\|A\phi_n-A\phi_m\|\le M\|\phi_n-\phi_m\|\rightarrow 0$$
Also konvergiert $A\phi_n\rightarrow B\phi\in\mathcal{H}$.\\
Sei $\psi_n\rightarrow\phi $ eine weitere Folge in $D(A)$. Dann gilt:
$$\|A\phi_n-A\psi_n\|\le M\|\phi_n-\psi_n\|\rightarrow 0$$
Also der Grenzwert $B\phi$ unabhängig von der approximierenden Folge.\\
$B$ ist linear. Sei $\phi_n,\psi_n\rightarrow\phi,\psi\in\mathcal{H}$ Folgen in $D(A)$. Dann gilt:
$$B(\phi+\lambda\psi)=\lim A(\phi_n+\lambda\psi_n)=\lim A\phi_n+\lambda \lim A\psi_n=B\phi+\lambda B\psi$$
$B$ ist eine Fortsetzung von $A$. Sei $\phi_n\rightarrow\phi\in D(A)$ eine Folge in $D(A)$. Dann gilt:
$$B\phi =\lim A\phi_n=A\phi $$
$B$ ist beschränkt. Sei $\phi\in D(A)$ mit $\|\phi\|\le 1$. Dann gilt $\|A\phi\|=\|B\phi\|$, also ist $\|A\|\le\|B\|$. \\
Sei $\phi_n\rightarrow\phi\in\mathcal{H}$ mit $\|\phi\|\le 1$ eine Folge in $D(A)$. Dann ist $\psi_n:=\frac{\|\phi\|\phi_n}{\|\phi_n\|+1/n}$ ebenfalls eine Folge in $D(A)$ die gegen $\phi$ konvergiert:
\begin{align*}
\|\psi_n-\phi\|&\le \|\psi_n-\phi_n\|+\|\phi_n-\phi\|=\|\phi_n\|\cdot\kla\left( {\frac{\|\phi\|}{\|\phi_n\|+1/n}-1} \right)+\|\phi_n-\phi\|\rightarrow 0
\end{align*}
Also folgt (Warum nicht einfach mit $\phi_n$???):
$$\|B\phi\|=\|\lim A\psi_n\|=\lim\|A\psi_n\|\le\sup \klammer\left\{ {\|A\chi\|\ |\ \chi\in D(A),\ \|\chi\|\le 1} \right\}$$
$B$ ist eindeutig. Sei $\phi_n\rightarrow\phi\in\mathcal{H}$ eine Folge in $D(A)$ und für $\psi\in D(A)$ sei $B_1\psi=A\psi=B_2\psi$:
\begin{align*}
\|B_1\phi-B_2\phi\|&\le\|B_1\phi-A\phi_n\|+\|A\phi_n-B_2\phi\|\le\|B_1\|\cdot\|\phi-\phi_n\|+\|B_2\|\cdot\|\phi-\phi_n\|\rightarrow 0
\end{align*}
\end{proof}
\section{Orts- und Impulsoperatoren}
a)
\begin{proof}
Es gilt $\phi_\alpha\in L^2$. Sei $R>0$ so, dass $\phi\vert_{[-R,R]}\equiv 0$. Dann gilt:
$$\int |\phi_\alpha|^2=\alpha\int|\phi(\alpha x)|^2\kom\overset{\substack{{y\,=\,\alpha x} \\ |}}{=}\intx\int_{-R}^{R}|\phi(y)|^2\,\mathrm{d}y<\infty $$
Allerdings sind $Q$ und $P$ unbeschränkt:
\begin{align*}
\|Q\phi_\alpha\|_2^2=\alpha\int |x\phi(\alpha x)|^2=\frac{1}{\alpha^2}\int |y\phi(y)|^2=\frac{1}{\alpha^2}\|Q\phi\|_2^2\rightarrow\infty \\
\|P\phi_\alpha\|_2^2=\alpha\int |\alpha\phi'(\alpha x)|^2=\alpha^2\int |y\phi'(y)|^2=\alpha^2\|P\phi\|_2^2\rightarrow\infty
\end{align*}
Die Einschränkung lässt $P$ unbeschränkt (Gleiches Argument, Einschränkung von $\phi_\alpha$ auf $(-1, 1)$). \\
$Q$ ist dann beschränkt:
$$\|Q\phi\|_2^2=\intx\int_{-1}^{1}|x\phi|^2\,\mathrm{d}x\le\intx\int_{-1}^{1}|\phi|^2\,\mathrm{d}x=\|\phi\|_2^2$$
\end{proof}
b)
\begin{proof}
Klar. Siehe Quantentheorie I.
\end{proof}
c)
\begin{proof}
Der Induktionsanfang entspricht den Voraussetzungen. Sei also $nB^{n-1}=AB^n-B^nA$:
\begin{align*}
(n+1)B^n&=B^n+nBB^{n-1}=B^n+BAB^n-B^{n+1}A\kom\overset{\substack{{[A,B]=I} \\ |}}{=}\\
&=B^n+ABB^n-IB^n-B^{n+1}A=AB^{n+1}-B^{n+1}A
\end{align*}
Weil $A,B$ beschränkt sind mit $\|A\|,\|B\|\le M<\infty$ folgt dann:
\begin{align*}
(n+1)\|B^n\|&\le \|AB^{n+1}\|+\|B^{n+1}A\|=2\|A\|\cdot\|B\|\cdot\|B^n\|\le 2M^2\|B^n\|
\end{align*}
Also $\exists N\in\mathbb{N}_0:\|B^N\|=0$, da sonst $(n+1)>2M^2$ für genügend große $n$ der Ungleichung widerspricht. \\
Andererseits gilt für $k\in\klammer\left\{ {1\dots N} \right\}$:
$$B^{N-k}=\frac{\klam\left[ {A, B^{N-k+1}} \right]}{N-k+1}$$
Daraus folgt dann $B^{n\le N}=0$. \\
Hieraus folgt der Widerspruch:
$$I=\klam\left[ {A,B} \right]=\klam\left[ {A, 0} \right]=0$$
\end{proof}
\section{Matrixexponentialfunktion}
a)
\begin{proof}
Es gilt $\|A^k\|\le\|A\|^k$. Sei $N\in\mathbb{N}$:
\begin{align*}
\|e^A\|&=\lim\Big\|\sum_k^N\frac{A^k}{k!}\Big\|\le\lim\sum_k^N\frac{\|A^k\|}{k!}\le\lim\sum_k^N\frac{\|A\|^k}{k!}= e^{\|A\|}
\end{align*}
\end{proof}
b)
\begin{proof}
Es gilt:
\begin{align*}
e^{A+B}&=\lim\sum_k^N\frac{(A+B)^k}{k!}=\lim\sum_k^N\sum_j^k\binom{j}{k}\frac{A^jB^{k-j}}{k!}=\lim\sum_k^N\sum_j^k\frac{A^j}{j!}\frac{B^{k-j}}{(k-j)!}=\\
&=\lim\kla\left( {\sum_n^N\frac{A^n}{n!}} \right)\cdot\kla\left( {\sum_m^N\frac{B^m}{m!}} \right)=e^Ae^B
\end{align*}
Betrachte:
\begin{align*}
I&=e^{A-A}=e^Ae^{-A}=e^{-A}e^A\ \ \Rightarrow\ \ e^{-A}=\kla\left( {e^A} \right)^{-1}
\end{align*}
\end{proof}
c)
\begin{proof}
Es gilt:
\begin{align*}
\frac{1}{h}\|e^{tA+hA}-e^{tA}-hAe^{tA}\|&=\frac{1}{h}\|e^{tA}\kla\left( {e^{hA}-1-hA} \right)\|=\frac{1}{h}\|e^{tA}\sum_{k=2}^\infty\frac{h^kA^k}{k!}\|\rightarrow 0
\end{align*}
\end{proof}
\end{document}